\section{Visual Inspection Results and Extrapolation}

\subsection{VI sample}

Visual inspection (VI) was carried out on the six tiles selected in
Section~\ref{sec:vi_selection} — three dark and three bright — comprising
a total of 391 outlier spectra (79 dark, 312 bright).
For each spectrum the inspector recorded whether an identifiable spectroscopic
problem was present, and whether the pipeline redshift quality flag was set
($\texttt{ZWARN} \neq 0$).
The per-tile results are given in Table~\ref{tab:vi_tiles}.

\begin{table}[h]
\centering
\caption{Per-tile VI results. $N$ is the total number of outlier spectra inspected,
$N_\mathrm{bad}$ the number with an identifiable VI problem,
$N_\mathrm{good}$ the number with no VI problem,
and $N_\mathrm{zw}$ the number with $\texttt{ZWARN}\neq 0$.}
\label{tab:vi_tiles}
\begin{tabular}{llrrrr}
\hline
Program & TILEID & $N$ & $N_\mathrm{bad}$ & $N_\mathrm{good}$ & $N_\mathrm{zw}$ \\
\hline
Dark   & 1406  &  10 &   5 &   5 &  1 \\
Dark   & 3473  &  30 &  23 &   7 &  2 \\
Dark   & 3833  &  39 &  35 &   4 & 10 \\
\hline
Bright & 25300 &  68 &  28 &  40 &  1 \\
Bright & 22175 & 101 &  45 &  56 &  1 \\
Bright & 22097 & 143 & 125 &  18 &  1 \\
\hline
\end{tabular}
\end{table}

\subsection{Aggregate VI fractions}

We compute 95\% Clopper--Pearson confidence intervals for two quantities:
the fraction of outliers with an identifiable VI problem ($f_\mathrm{bad}$)
and the fraction with $\texttt{ZWARN}\neq 0$ ($f_\mathrm{zw}$).
The Clopper--Pearson interval for $k$ events in $n$ trials at confidence level
$1-\alpha$ is
\begin{equation}
    \left[
        B\!\left(\tfrac{\alpha}{2};\,k,\,n-k+1\right),\;
        B\!\left(1-\tfrac{\alpha}{2};\,k+1,\,n-k\right)
    \right],
\end{equation}
where $B(p;\,a,\,b)$ denotes the $p$-th quantile of the Beta distribution.

Summing over all tiles within each program the aggregated results are
shown in Table~\ref{tab:vi_agg}.

\begin{table}[h]
\centering
\caption{Aggregated VI fractions with 95\% Clopper--Pearson intervals.}
\label{tab:vi_agg}
\begin{tabular}{lrrrr}
\hline
Program & $N$ & $f_\mathrm{bad}$ [\%] & $N_\mathrm{good}$ & $f_\mathrm{zw}$ [\%] \\
\hline
Dark   &  79 & $79.7\ [69.2,\ 88.0]$ & 16  & $16.5\ [9.1,\ 26.5]$ \\
Bright & 312 & $63.5\ [57.9,\ 68.8]$ & 114 & $1.0\ [0.2,\ 2.8]$   \\
\hline
All    & 391 & $66.8\ [61.8,\ 71.4]$ & 130 & $4.1\ [2.4,\ 6.6]$   \\
\hline
\end{tabular}
\end{table}

A striking result is that $\texttt{ZWARN}$ catches very few of the
outliers identified by UMAP: only $4.1\%$ of all inspected spectra
have $\texttt{ZWARN}\neq 0$, with the rate being particularly low
in the bright program ($1.0\%$).
This confirms that the UMAP outliers represent a population of
spectroscopically problematic spectra that are largely invisible to
the standard pipeline quality flag.

\subsection{Extrapolation to the full Loa main survey}

We extrapolate the VI good fractions ($1 - f_\mathrm{bad}$) measured
in the sample to the full set of dark and bright main-survey outliers
in Loa (191{,}411 dark and 463{,}543 bright spectra; see
Table~\ref{tab:loa_summary}).
The estimated number of outliers without an identifiable VI problem is
\begin{equation}
    \hat{N}_\mathrm{good} = (1 - \hat{f}_\mathrm{bad})\times N_\mathrm{total},
\end{equation}
where $\hat{f}_\mathrm{bad}$ is the sample fraction and the uncertainty
is propagated through the Clopper--Pearson interval.
The results are given in Table~\ref{tab:vi_extrap}.

\begin{table}[h]
\centering
\caption{Estimated number of no-VI-problem outliers in the full Loa main survey,
extrapolated from the VI sample fractions.
Brackets give the 95\% Clopper--Pearson interval.}
\label{tab:vi_extrap}
\begin{tabular}{lrrr}
\hline
Program & Good fraction [\%] & $N_\mathrm{total}$ & $\hat{N}_\mathrm{good}$ \\
\hline
Dark   & $20.3\ [12.0,\ 30.8]$ & 191{,}411 & $38{,}767\ [23{,}051,\ 58{,}946]$ \\
Bright & $36.5\ [31.2,\ 42.1]$ & 463{,}543 & $169{,}371\ [144{,}558,\ 195{,}382]$ \\
\hline
Combined & $33.2\ [28.6,\ 38.2]$ & 654{,}954 & $217{,}760\ [187{,}271,\ 249{,}924]$ \\
\hline
\end{tabular}
\end{table}

Combining both programs, we estimate that approximately
$\mathbf{218{,}000}$ outliers in the Loa dark and bright main survey
have no identifiable VI problem, with a 95\% confidence interval of
$[187{,}000,\ 250{,}000]$.
These spectra are genuine UMAP outliers that are not flagged by the
standard pipeline and warrant further investigation.
\end{document}
